\section{Two positive-domain improvements by combining the wave-10 mechanisms}
\label{sec:positive-improvements}

The integer-domain construction in the preceding sections is not the only
useful endpoint.  The same review yields two improvements to the positive
frontier.  They are separate, fully specified polynomials with the authentic
compact44 coefficient tuple.  Their contribution is a precise composition
of the mechanisms in sources~69 and~70, rather than a claim that the generic
norm or descent methods are new.

\begin{center}
\begin{tabular}{lrrr}
\hline
Construction & Positive variables & Exact degree & Operations\\
\hline
Source 70 & 11 & 325 & 113\\
Guard-free strong ordinate & 11 & 277 & 111\\
Source 69 & 12 & 82 & 102\\
Adaptive main carrier & 12 & 78 & 103\\
\hline
\end{tabular}
\end{center}

The operation convention counts every live binary addition, subtraction,
and multiplication on witnesses, including squaring and scalar
multiplication. Fixed integer leaves and register reuse are free. These are
arithmetic-circuit counts, not bit-complexity estimates.

\subsection{An adjacent-coefficient descent}

We use the following arithmetic lemma from source~69\cite{source69}.
Its role here is to remove a hypothesis that source~70 enforced by an
additional factor in its strong Pell ordinate.

\begin{lemma}[Source 69: rational ordinate descent]
\label{lem:adjacent-descent}
Let $Q>1$ and $V$ be integers. If $y\in\mathbb Q$ and
\[
 QV^2-(Q-1)y^2=1,
\]
then $y\in\Z$.
\end{lemma}
\begin{proof}
Write $y=a/m$ in lowest terms, with $m>0$. The equation implies
$m^2\mid Q-1$, and excludes $V=0$ and $y=0$. Change signs to make both
positive. If $V=1$ then $y=1$. Otherwise define
\[
 V'=(2Q-1)V-2(Q-1)y,\qquad
 y'=(2Q-1)y-2QV.
\]
The first is integral because $m\mid Q-1$. Direct expansion preserves the
norm. The reduced denominator of $y'$ is exactly $m$, since
$2Q-1\equiv1\pmod m$. Moreover $y>V$ and $QV>(Q-1)y$, and hence
$-V<V'<V$. The norm excludes $V'=0$. Strong induction on the positive
integer $V$, after taking absolute values, reaches the case $V=1$
without changing the denominator. Thus $m=1$.
\end{proof}

The exact polynomial norm identity in this proof is checked by the supplied
twelve-variable verifier. No finite enumeration is being used in place of
the induction.

\subsection{Eleven positive unknowns in degree 277}

Use the eleven inputs and every register of source~70's eleven-positive
construction\cite{source70}, with exactly one change:
\[
 \boxed{I=ic^2\quad\hbox{in place of}\quad I=iq(q-1)c^2.}
\]
The inputs are
\[
 (j,F,\alpha,i,T,X,Y,n,\eta,\zeta,\gamma)\in\Pos^{11}.
\]
The fixed numerals are $\beta,K_0,M_C,M_F$, with
$\beta=2^{590069}-1$ and the already-shifted $M_F$.
All abbreviations below are polynomial registers, not extra unknowns.
For completeness the definition is
\begin{align*}
 J&=3j+1,&q&=\beta J+1,&C&=q-F-\alpha,\\
 E&=XY,&k&=\eta+\zeta,&c&=kY+\eta,\\
 a&=E+Y,&H&=4a+3,&\Delta&=a^2+H,\\
 \ell&=EkY,&I&=ic^2,&M&=q^3(q-1),\\
 R&=((q-1)(q-F)+\alpha)(q^2-1)+(M_C+qM_F)J,\\
 R_+&=q^3(q+F+\alpha),&
 B_t&=(qK_0+X)C-qF,&D_\tau&=\ell(\ell+k)+1,\\
 D_f&=1+\Delta I^2,&Q&=\Delta(D_f-1),&\delta&=Q-1,\\
 V&=D_f(cT-R)-c,&L&=3D_f(cT+R_+),&W&=1+(q-1)L,\\
 Z&=M(\ell+k+n+L)+(q-1)Y-q^3WB_t,\\
 S&=4qY,&U&=qS,&v&=k-R-1,\\
 D_*&=U(X+ac)+H(U\gamma+SX-v),\\
 A_1&=D_*^2-U^2(1+\Delta c^2).
\end{align*}
Now put
\begin{align*}
 A&=M^2D_\tau,&B_f&=M^2D_f,&C_y&=M^2(QV^2-1),\\
 N_c&=\delta^2\bigl[(Z^2-A-B_f)^2-4AB_f\bigr]
       -2\delta C_y(Z^2+A+B_f)+C_y^2,\\
 N_8&=N_c^2-64\delta^3AB_fC_yZ^2.
\end{align*}
The required equation is
\begin{equation}
 \boxed{P_{11}:=(1+N_8)(1+A_1)-1=0.}
 \label{eq:positive277}
\end{equation}
The exact division-free definition is also supplied as
\texttt{positive/eleven/p11\_positive\_degree277.json}.

\begin{prop}
ZFC proves that \eqref{eq:positive277} has no zero in $\Pos^{11}$ if and
only if $\operatorname{Con}(\mathrm{ZFC})$. Its exact total degree is 277
and its emitted evaluator has 111 live binary operations.
\end{prop}
\begin{proof}
We give the complete change to the upstream proof; its normalized Pell
and native compiler interfaces remain the ones in sources~61 and~70
\cite{source61,source70}. The source~70 inequalities use only $I\ge c^2$,
which still holds. In particular, writing
\[
 D_y=\frac{QV^2-1}{Q-1},
\]
the three positive rational radicands $D_\tau,D_f,D_y$ are pairwise
distinct, and no signed sum of their square roots is zero. The octic norm
therefore recovers rational signed roots $\tau_*,f_*,y_*$ with squares
$D_\tau,D_f,D_y$. The first two are integers because their squares are
integers. Lemma~\ref{lem:adjacent-descent} makes the third an integer.
Here $Q,V$ are integers and $Q>1$ before any equation is imposed.

This replaces the old argument that the denominator of $y_*$ divides $M$
and is coprime to $Q-1$. That argument was the only reason for the extra
$q(q-1)$ in $I$. The unchanged recovery formulas are
\[
 f_{\rm old}=|f_*|,\quad y_{\rm old}=|y_*|,\quad
 u=y_{\rm old}-V>0,\quad r=n+L-f_*-y_*>0.
\]
The signed roots must be used in the formula for $r$. The source proof's
bound $L>|f_*|+|y_*|$ makes it positive. The strong ordinate $I=ic^2$
has exactly the divisibility needed by the normalized index theorem.
Consequently the original four-residual positive source interface is
recovered, and its calibration yields a finite syntactic ZFC refutation.

Conversely any positive zero of source~70 maps directly to a zero here:
replace its input $i$ by $q(q-1)i$ and leave every other input unchanged.
This preserves $I$ and every downstream register. Hence source~70's
completeness supplies completeness here with no new index restriction.
The unit-product finalizer is unchanged: both $N_8$ and $A_1$ are squares
modulo 4, so neither can equal $-2$.

For exact degree, the changed register degrees are
\[
\begin{array}{c|rrrrrrrrr}
\text{register}&I&D_f&Q&L&W&Z&N_c&N_8&A_1\\\hline
\text{degree}&5&14&18&18&19&24&132&264&13.
\end{array}
\]
Writing $q_0=3\beta j$, $k_0=\eta+\zeta$ and
$B_0=(K_0q_0+X)(q_0-F-\alpha)-q_0F$, the leading forms are
\begin{align*}
 Q_0&=i^2k_0^4X^4Y^8,\\
 Z_0&=-3i^2q_0^7k_0^4X^2Y^6(q_0+F+\alpha)B_0,\\
 (A_1)_{13}&=32q_0^2k_0X^2Y^4
      \bigl(4q_0^2Y\gamma+q_0^3(q_0-F)\bigr).
\end{align*}
The unique top term of $N_8$ is $Q_0^4Z_0^8$. Therefore the coefficient
of
\[
 j^{60}i^{24}\eta^{49}X^{42}Y^{85}\alpha^{16}\gamma
\]
in $P_{11}$ is $2^7 3^8(3\beta)^{60}>0$. The monomial has degree 277.
The exact checker verifies these identities and the literal graph count:
65 multiplications, 29 additions and 17 subtractions.
\end{proof}

\subsection{Twelve positive unknowns in degree 78}

For a second endpoint, keep all source~69 twelve-variable registers
\cite{source69}, except replace its main Pell carrier by
\begin{align*}
 S&=4qY, & U&=qS=4q^2Y,\\
 D_*&=U(X+ac)+H\bigl(U\gamma+SX-(k-R-1)\bigr),\\
 A_1&=D_*^2-U^2(1+\Delta c^2).
\end{align*}
The source's first carrier $A_0$, two-root resultant $N$, and finalizer
are unchanged:
\begin{equation}
 \boxed{P_{12}=(1+A_0)(1+A_1)(1+N)-1=0.}
 \label{eq:positive78}
\end{equation}
Its full literal definition is
\texttt{positive/twelve/equation.json}; the twelve positive inputs are
$j,F,\alpha,t,T,X,Y,r,\eta,\zeta,n,\gamma$.

To make the unchanged part explicit, use $J=6j+1$, $q=\beta J+1$,
$C=q-F-\alpha$, and the same $E,k,c,a,H,\Delta,\ell,R,B_t$ formulas
as above. Put
\begin{align*}
 W_0&=1+(q-1)k,& M&=q^3(q-1),\\
 Z_0&=M(\ell+k+r)+(q-1)Y-q^3W_0B_t,\\
 A_0&=Z_0^2-M^2\bigl(\ell(\ell+k)+1\bigr),\\
 f&=t+1,& \rho&=f^2-1,& d&=\Delta c^4,\\
 Q&=\Delta\rho,& V&=f^2(cT-R)-c,\\
 h(z)&=(1-Q)z^2+2(1-Q)Vz+V^2-1,\\
 A_N&=dh(n)+(1-Q)\rho,& B_N&=2(1-Q)(n+V),\\
 N&=A_N^2-dB_N^2\rho.
\end{align*}
All coefficients are fixed integers and all powers have fixed exponents.

\begin{prop}
ZFC proves that \eqref{eq:positive78} has no zero in $\Pos^{12}$ if and
only if $\operatorname{Con}(\mathrm{ZFC})$. Its exact total degree is 78
and its emitted evaluator has 103 live binary operations.
\end{prop}
\begin{proof}
The unit-product finalizer still works: $q$ is even, $4\mid U$, and all
three residuals are squares modulo 4. First $A_0=0$ gives $q^3\mid Y$
and $C>0$, hence $q\ge16$, $q\equiv2\pmod3$ and
$3q<R<q^4-2$. The new $D_*$ is positive, since
$D_*>Uac-Hk>ak(UY-5)>0$. Thus $A_1=0$ recovers the positive main
Pell root. The first and main roots give $c>2R$ before the strong or
auxiliary equations are used. Consequently $V>1$, and the unchanged
resultant and Lemma~\ref{lem:adjacent-descent} recover positive integral
strong and auxiliary coordinates. The normalized auxiliary-index theorem
then gives $p=R$.

For the altered carrier set $\mu=S/\gcd(S,H)$. Since $q^3\mid Y$,
$q\equiv2\pmod3$ and $H=4Y(X+1)+3$, one has $\gcd(S,H)\le3$.
Thus $\mu\ge4qY/3>q^4>R+1$. Both $S$ and $U$ divide $4XY^2$
without assuming $q\mid X$. Carrier divisibility gives
$\mu\mid k-R-1$, and the first Pell recurrence gives
$k\equiv2n_{\rm first}\pmod{4XY^2}$. With
$1\le n_{\rm first}<p=R<\mu$, these imply
$2n_{\rm first}=R+1$. Therefore $v=k-R-1$ is a positive multiple of
$4XY^2$, hence of $U$. Returning to the carrier gives $q\mid HX$,
and $\gcd(q,H)=1$ gives $q\mid X$.

Writing $w=X/q$, the restored source gap is
$\gamma_{\rm old}=\gamma+w-v/U$. It is integral and positive because
the main root is greater than $(a+1)c$ and $c>X$.
This recovers the source~69 interface. Conversely its canonical witnesses
lift by $\gamma=\gamma_{\rm old}+v/U-w>0$, using
$v/U\ge XY/q^2=wY/q\ge wq^2>w$. No new canonical restriction is
introduced.

The exact residual degrees are $15,13,50$. Their product has degree 78:
its coefficient at
\[
 j^{18}F^2\eta^{11}X^{12}Y^{23}t^8n^2\gamma\alpha
\]
is $1024(6\beta)^{18}>0$. The supplied checker verifies the complete
carrier expansions, resultant identities and this coefficient. It counts
58 multiplications, 26 additions and 19 subtractions. Thus this endpoint
trades one operation for four degrees compared with source~69.
\end{proof}

These two improvements preserve the native compiler and syntactic
consistency calibration. Their exact algebra is independently replayable,
but they are not presented as new end-to-end proof-assistant certificates.
